% Article court : la note sur les nombres triangulaires de NAF 4073, f. 32v.
% Compiler : tectonic -X compile docs/naf-4073/article.tex
%
% Destiné à Hexagon (hexagonmath.org) : sources TeX seules, sans code Lean.
% À FAIRE avant le dépôt : relire la note sur l'image de Gallica (vue 42) et
% la comparer à la fig. 2 d'Alexanderson 2012. Après relecture, la section
% « Sources » peut dire que les lectures ont été contrôlées.
\documentclass[11pt,a4paper]{article}
\usepackage{iftex}
\ifPDFTeX
  \usepackage[T1]{fontenc}
  \usepackage[utf8]{inputenc}
  \usepackage{lmodern}
\else
  \usepackage{fontspec}
\fi
\usepackage[french]{babel}
\usepackage{amsmath,amssymb,amsthm}
\usepackage{graphicx}
\usepackage[margin=2.8cm]{geometry}
\usepackage[hidelinks]{hyperref}

\newtheorem{proposition}{Proposition}
\theoremstyle{remark}
\newtheorem*{remarque}{Remarque}

\title{«~Le triangulaire d'un carré~» : une note inédite\\
attribuée à Sophie Germain (BnF, NAF 4073)}
\author{Michel Hua\\\small Chercheur indépendant}
\date{Novembre 2026}

\begin{document}
\maketitle

\begin{abstract}
Le volume NAF 4073 de la Bibliothèque nationale de France réunit des lettres
adressées à Sophie Germain ou écrites par elle. Au verso du feuillet~32, qui
a servi de couverture à une lettre non signée du 13~mars 1814, une note
d'une autre main que celle de la lettre énonce : «~le triangulaire d'un
quarré est la somme des quarrés des deux Triangulaires consécutifs~», et
l'écrit
$\big(\frac{n^2-n}{2}\big)^2+\big(\frac{n^2+n}{2}\big)^2=\frac{n^4+n^2}{2}$.
C'est la seule ligne de mathématiques du volume ; ni les catalogues ni les
études que j'ai consultées n'en font mention. Je propose de l'attribuer à
Germain. Un manuscrit qui lui est attribué, daté «~vers 1797~» par le
libraire qui l'a vendu et reproduit par Alexanderson \cite{alexanderson},
donne la même identité, dans la même notation, après les mots «~J'ai aussi
remarqué que~», et l'écrit en symboles exactement comme la phrase de
NAF 4073 la dit en mots. L'identité n'est pas nouvelle ; la note l'est.
\end{abstract}

Au verso d'une enveloppe adressée à Sophie Germain en 1814, une main
rapide a écrit une formule et la phrase qui la lit : le triangulaire d'un
carré est la somme des carrés de deux triangulaires consécutifs. Aucun
catalogue ni aucune étude ne signale cette note. Je la décris
(section~\ref{sec:feuillet}), j'en donne le sens et la preuve
(section~\ref{sec:identite}), puis j'expose pourquoi je la crois de Germain
et ce que cela permet de dire de sa date (section~\ref{sec:main}). Ce qui
reste ouvert est rassemblé à la fin.

\section{Le feuillet}\label{sec:feuillet}

NAF 4073 vient des papiers de Guillaume Libri, et les catalogues n'y voient
que des lettres : «~Lettres adressées à Sophie Germain ou écrites par elle~»
pour Delisle \cite[p.~183]{delisle} ; «~Lettres de Delambre, Joseph Fourier,
Sophie Germain, Lagrange, Lalande et G.~Libri~» pour la notice actuelle de la
BnF. Le volume est numérisé dans Gallica \cite{gallica}.

Le feuillet~32 enveloppait une lettre non signée du 13~mars 1814, que
Charles Henry a publiée en entier \cite[n\textsuperscript{o}~XVII,
p.~637--638]{henry}. Le verso (vue~42 de Gallica) porte l'adresse,
«~Mademoiselle Germain, rue Sainte-Croix-de-la-Bretonnerie~», à Paris, dont
le numéro a été recouvert après coup par deux larges boucles d'encre, et un
cachet postal. En haut du feuillet, tête-bêche par rapport à l'adresse, une
autre main a écrit une note (figure~\ref{fig:f32v}). La voici telle que la transcrit l'édition, avec
son apparat ; les crochets marquent la restitution, les appels renvoient aux
notes de l'éditeur.

\begin{center}
\fbox{\begin{minipage}{0.9\textwidth}\small
\medskip
\centerline{\textsc{NAF 4073, f.~32v (vue 42)}}
\[
\left(\frac{(n^2-n)}{2}\right)^2 + \left(\frac{(n^2+n)}{2}\right)^2
= \frac{n^4+n^2}{2}
\]
\smallskip
$1 + 64 \quad 189$\,\textsuperscript{a}

\smallskip
le triangulaire d'un quarré est la somme des quarrés des deux Triangulaires
consécu[tifs]\,\textsuperscript{b,\,c}

\medskip
\hrule
\medskip
\footnotesize
\textsuperscript{a}\,Les nombres sont écrits sur une seule ligne ; un trait
horizontal continu part du signe plus et traverse $64$ jusqu'au premier
chiffre de $189$, qu'une première lecture avait pris pour un second signe
plus suivi de $49$. Relu $189$ sur l'agrandissement de la vue par un humain
(Michel Hua), le 25~novembre 2026. Le trait peut être une rature. Rien ne
rattache explicitement cette ligne à l'égalité qui la précède.

\textsuperscript{b}\,Le dernier mot est coupé par la reliure au bord du
feuillet : les lettres «~consécu~» sont entières, la suite est prise dans le
fond du cahier.

\textsuperscript{c}\,Cette note est écrite à l'encre, tête-bêche par rapport
à l'adresse, en haut du feuillet ; la main n'est pas celle de la lettre. Elle
est plus rapide et plus anguleuse, et rappelle celle de la lettre du
18~avril 1831 aux vues~27 à~29, sans que quinze mots suffisent à
l'établir : la présente lecture ne tranche pas.
\medskip
\end{minipage}}
\end{center}

\begin{figure}[ht]
\centering
\includegraphics[width=0.72\textwidth]{naf-4073-f32v.png}
\caption{BnF, NAF 4073, f.~32v (vue~42), partie haute : la formule, la ligne
$1 + 64 \quad 189$ et, à droite, la phrase ; en bas à gauche, l'adresse.
Source : gallica.bnf.fr / Bibliothèque nationale de France,
\url{https://gallica.bnf.fr/ark:/12148/btv1b10089787b/f42.item}.
Reproduction exclue de la licence de l'article.}
\label{fig:f32v}
\end{figure}

En graphie moderne : «~le triangulaire d'un carré est la somme des carrés
des deux triangulaires consécutifs~». Sur le feuillet, la formule est en haut et la phrase plus bas à droite, en
cinq lignes courtes ; la ligne de nombres est à gauche, sous la formule. La
note~c est celle de la transcription, écrite avant que je connaisse le
manuscrit publié par Alexanderson ; la section~\ref{sec:main} dit pourquoi
je vais aujourd'hui plus loin.

Henry ne dit pas un mot de cette note, alors qu'il imprime la lettre qu'elle
recouvre. Je ne l'ai trouvée ni dans les catalogues \cite{delisle,omont}, ni
dans l'étude des lettres de Germain conservées à Florence \cite{delcentina},
ni dans les travaux sur ses manuscrits de théorie des nombres
\cite{laubenbacher-pengelley}. L'inventaire de sa correspondance que donnent
Del~Centina et Fiocca \cite[annexe]{delcentina-fiocca} recense la lettre du
13~mars 1814 et ignore sa couverture. Le corps de leur chapitre, en accès
payant, je ne l'ai pas lu.

\section{L'identité}\label{sec:identite}

Notons $T_k = 1 + 2 + \cdots + k = k(k+1)/2$ le $k$-ième nombre
triangulaire, la somme des entiers de $1$ à $k$, ainsi nommé parce que
$T_k$ points se rangent en triangle de côté $k$ : $T_1 = 1$, $T_2 = 3$,
$T_3 = 6$, $T_4 = 10$\dots{} «~Le triangulaire d'un carré~» est
$T_{n^2}$, le triangulaire d'indice $n^2$ ; «~le carré d'un triangulaire~»
est $(T_k)^2$.

La formule $1 + 2 + \cdots + k = k(k+1)/2$ est celle de l'anecdote de Gauss
écolier : prié de faire la somme des entiers de $1$ à $100$, il l'aurait
donnée sur-le-champ, $5050$, en associant les termes deux à deux,
$1 + 100 = 2 + 99 = \cdots = 101$, cinquante fois. En général, écrire la
somme à l'endroit et à l'envers donne $k$ paires valant chacune $k+1$, d'où
$2T_k = k(k+1)$. L'histoire a enflé avec les récits ; le plus ancien, celui
de Sartorius von Waltershausen en 1856, ne donne ni la suite ni le nombre
$100$ \cite{hayes}. La formule, elle, est bien plus vieille que Gauss.

Sur le feuillet, les deux nombres élevés au carré,
$(n^2-n)/2$ et $(n^2+n)/2$, sont $T_{n-1}$ et $T_n$, et le second membre,
$(n^4+n^2)/2 = n^2(n^2+1)/2$, est $T_{n^2}$. La note dit donc
\begin{equation}\label{eq:note}
(T_{n-1})^{2} + (T_n)^{2} = T_{n^2}.
\end{equation}
Pour $n = 2$, $1^2 + 3^2 = 10 = T_4$ ; pour $n = 3$, $3^2 + 6^2 = 45 = T_9$ ;
pour $n = 4$, $6^2 + 10^2 = 136 = T_{16}$. La suite $1, 10, 45, 136,
325, \dots$ des $T_{n^2}$ est la suite A037270 de l'OEIS, où l'identité
figure parmi les formules depuis 2013, sans référence historique
\cite{oeis}. La phrase française dit cela, et
un peu plus : «~les deux Triangulaires consécutifs~» ne sont pas n'importe
quels triangulaires qui se suivent, ce sont les deux dont le carré $n^2$ est
la somme, puisque $T_{n-1} + T_n = n^2$.

\begin{proposition}\label{prop:note}
Pour tout entier $n \geq 1$, $(T_{n-1})^2 + (T_n)^2 = T_{n^2}$.
\end{proposition}

\begin{proof}
Je suis le calcul du manuscrit reproduit par Alexanderson
(figure~\ref{fig:quaritch}). On développe les deux carrés :
\[
\left(\frac{n^2+n}{2}\right)^2 = \frac{n^4 + 2n^3 + n^2}{4},
\qquad
\left(\frac{n^2-n}{2}\right)^2 = \frac{n^4 - 2n^3 + n^2}{4}.
\]
Dans la somme, les termes $\pm 2n^3$ se détruisent :
\[
(T_n)^2 + (T_{n-1})^2 = \frac{2n^4 + 2n^2}{4} = \frac{n^4 + n^2}{2}
= \frac{n^2(n^2+1)}{2} = T_{n^2}.
\]
\end{proof}

Le même développement donne les deux résultats du manuscrit. Dans la
différence, ce sont les termes $n^4$ et $n^2$ qui se détruisent, et il reste
$(T_n)^2 - (T_{n-1})^2 = 4n^3/4 = n^3$, d'où la proposition de Lagrange
en ajoutant ces égalités de $1$ à $n$ ; dans la somme, ce sont les cubes qui
s'en vont, et il reste $T_{n^2}$. On peut aussi raisonner sans développer :
avec $a = T_{n-1}$ et $b = T_n$, on a $a + b = n^2$ et $b - a = n$, donc
$a^2 + b^2 = \frac12\big((a+b)^2 + (b-a)^2\big) = \frac12(n^4+n^2)$. La
phrase du feuillet suit plutôt cet ordre-là : un carré, puis les deux
triangulaires dont il est la somme et dont sa racine est la différence.
Rien de tout cela ne dépasse ce qu'on savait faire en 1797.

L'identité est aussi le cas $m = n$ d'une formule plus générale, facile à
vérifier et dont je ne connais pas l'origine :
\[
T_{mn} = T_m\,T_n + T_{m-1}\,T_{n-1},
\]
car $4\,(T_m T_n + T_{m-1} T_{n-1}) = mn\big[(m+1)(n+1) + (m-1)(n-1)\big]
= mn\,(2mn + 2) = 4\,T_{mn}$. Les termes croisés $\pm(m+n)$ s'y détruisent,
comme les $\pm 2n^3$ dans le calcul de Germain.

Si l'on joint \eqref{eq:note} à l'identité de Nicomaque,
$1^3 + \cdots + n^3 = (T_n)^2$, on a $(T_{n-1})^2 = (T_n)^2 - n^3$, et une
forme que le feuillet ne donne pas :

\begin{proposition}
Pour tout $n \geq 1$, $T_{n^2} = 2\,(T_n)^2 - n^3$ : le triangulaire d'un
carré vaut deux fois la somme des $n$ premiers cubes, moins le dernier. Pour
$n = 3$ : $T_9 = 2 \cdot 36 - 27 = 45$.
\end{proposition}

\begin{remarque}[La ligne $1 + 64 \quad 189$]
La ligne n'instancie pas l'identité : $1 + 64 = 65$, et $189$ n'est ni
triangulaire ($T_{18} = 171$, $T_{19} = 190$) ni carré. Mais $1 = 1^3$,
$64 = 4^3$, et
\[
189 = 4^3 + 5^3 = T_{25} - T_{16} = T_{5^2} - T_{4^2}.
\]
C'est le pas de la suite $T_{n^2}$ entre $n = 4$ et $n = 5$, et ce pas vaut
toujours $n^3 + (n-1)^3$ : par \eqref{eq:note},
$T_{n^2} - T_{(n-1)^2} = (T_n)^2 - (T_{n-2})^2$, et chaque
$(T_k)^2 - (T_{k-1})^2$ vaut $k^3$ (c'est la différence de Germain,
figure~\ref{fig:quaritch}). C'est aussi la somme des $2n-1$ entiers de
$(n-1)^2 + 1$ à $n^2$, dont la moyenne est $n^2 - n + 1$ : pour $n = 5$,
$17 + 18 + \cdots + 25 = 9 \times 21 = 189$. Plus généralement, pour
$m < n$, $(T_n)^2 - (T_m)^2 = (m+1)^3 + \cdots + n^3$. La ligne serait alors un essai qui relie la note
aux cubes de la question de Lagrange. Ce n'est qu'une hypothèse : le feuillet
ne pose pas ce calcul, et un nombre lu après coup se prête toujours à
quelque propriété.
\end{remarque}

\section{La main et la date}\label{sec:main}

C'est Georg Essl qui a fait le rapprochement, en réponse à la question que
j'avais posée sur History of Science and Mathematics \cite{essl}. En 2012,
G.~L.~Alexanderson a reproduit un manuscrit non signé, vendu par la librairie
Bernard Quaritch dans les années 1990, daté par elle «~vers 1797~» et
attribué à Germain ; il l'a fait reproduire en couverture du
\emph{Bulletin of the AMS} d'avril 2012 \cite{alexanderson}. La page commence
par «~La Grange m'ayant proposé de démontrer / Que le quarré d'un
triangulaire quelconque est égal à la somme des cubes de tous les nombres
naturels pris depuis 1 jusqu'à $n$ inclusivement~». Avec $T_n = 1 + 2 +
\cdots + n$, c'est l'identité de Nicomaque,
$(1 + 2 + \cdots + n)^2 = 1^3 + 2^3 + \cdots + n^3$. La preuve tient dans
une différence :
$\big(\frac{n^2+n}{2}\big)^2 - \big(\frac{n^2-n}{2}\big)^2 = n^3$,
c'est-à-dire $(T_n)^2 - (T_{n-1})^2 = n^3$ ; en ajoutant ces égalités de
$1$ à $n$, les termes se compensent deux à deux et il reste
$(T_n)^2 = 1^3 + \cdots + n^3$. Sous un trait vient la somme des deux mêmes
carrés : «~J'ai aussi remarqué que~»
$\big(\frac{n^2+n}{2}\big)^2 + \big(\frac{n^2-n}{2}\big)^2 = \frac{n^4+n^2}{2}$,
«~donc~» $\Delta^a + \Delta^{a-1} = \square^a$ et
\[
\square^{\Delta^a} + \square^{\Delta^{a-1}} = \Delta^{\square^a}
\]
\cite[p.~330]{alexanderson} (figure~\ref{fig:quaritch}). Germain note $\Delta^a$ le triangulaire de $a$,
notre $T_a$, et $\square^a$ le carré de $a$ ; l'exposant est un indice, non
une puissance. La première ligne dit donc $T_a + T_{a-1} = a^2$, et la
seconde $(T_a)^2 + (T_{a-1})^2 = T_{a^2}$, qui est \eqref{eq:note} avec $a$
pour $n$. Lue à voix haute, cette seconde ligne est la phrase de NAF 4073 :
le carré du triangulaire de $a$ et celui du triangulaire de $a-1$ font le
triangulaire du carré de $a$. Alexanderson ne parle pas de NAF 4073.

\begin{figure}[ht]
\centering
\includegraphics[width=\textwidth]{quaritch-germain.png}
\caption{Le manuscrit attribué à Germain, vendu par Bernard Quaritch et daté
par lui «~vers 1797~», tel que le reproduit Alexanderson
\cite[fig.~2, p.~330]{alexanderson}. En haut, la question de Lagrange et la
différence qui donne $n^3$ ; sous le trait, «~J'ai aussi remarqué que~», la
somme, puis les mêmes égalités écrites avec $\Delta$ et $\square$.}
\label{fig:quaritch}
\end{figure}

Trois sortes d'indices portent l'attribution, et ils n'ont pas le même
poids. Le plus fort est textuel. Les deux témoins écrivent la formule de la
même façon : mêmes termes $(n^2 \pm n)/2$ entre parenthèses, l'exposant hors
de la parenthèse, même second membre $(n^4+n^2)/2$, et non $n^2(n^2+1)/2$,
qui viendrait aussi naturellement ; et la phrase de NAF 4073 dit en mots ce
que le manuscrit Quaritch dit en symboles. Le deuxième est matériel : la
note est écrite sur une lettre adressée à Germain, chez elle, et conservée
dans ses papiers, passés ensuite par Libri. Le troisième, le plus faible,
est paléographique. Autant que j'en puisse juger sur les reproductions,
l'écriture est compatible avec celle de la figure~2 d'Alexanderson et avec
la lettre de Germain à Libri du 18~avril 1831 reliée dans le même volume
(vues~27 à 29) : même inclinaison, même \emph{q} de «~quarré~», même
liaison dans «~est~». Je ne suis pas paléographe. L'ensemble rend
l'attribution à Germain très probable ; il ne la démontre pas.

Le support date la note par en bas : elle est postérieure au 13~mars 1814,
puisque le feuillet avait d'abord servi d'enveloppe. Le manuscrit Quaritch
n'est daté que par le libraire. Si les deux témoins sont bien de Germain, la
note de NAF 4073 montre qu'elle est revenue à cette identité au moins
seize ans après la date qu'on donne à la première rédaction ; pourquoi, je
l'ignore.

L'identité est classique, et l'on savait par Alexanderson que Germain la
connaissait. Ce que le feuillet ajoute tient au manuscrit : un volume que
tous les catalogues décrivent comme une correspondance contient une ligne de
théorie des nombres, très probablement de sa main, que personne n'avait relevée.

\section*{Ce qui reste ouvert}

Une expertise paléographique de la note, à partir du fac-similé de Gallica
et de la figure~2 d'Alexanderson, confirmerait ou non l'attribution. Le fonds
Germain de la Biblioteca Moreniana, à Florence, plus de deux cents feuillets
\cite{delcentina}, est l'endroit où chercher une troisième copie, ou la
raison de celle-ci. Le corps du chapitre de Del~Centina et Fiocca
\cite{delcentina-fiocca}, que je n'ai pas lu, mentionne peut-être la note.
Restent enfin la ligne $1 + 64 \quad 189$, dont le sens n'est qu'une
hypothèse, et un
détail de lecture : sur le fac-similé, le «~s~» final de «~Triangulaires~»
est douteux.

\section*{Sources et vérification}

La transcription et la lecture modernisée du volume sont publiées sous
licence CC0, dans le répertoire \texttt{transcripts/naf-4073} du dépôt
\url{https://github.com/Commutator-IO/germain-project} ; la note est à la
vue~42, dans \texttt{batch-03.fr.tex}. La transcription a été faite avec un
modèle de langue lisant les images de Gallica ; les lectures citées ici
doivent être contrôlées sur le fac-similé.

Les deux propositions et les constats arithmétiques de la remarque sont
démontrés en Lean~4 avec mathlib, dans le fichier
\href{https://github.com/Commutator-IO/germain-project/blob/main/lean/Germain/Naf4073.lean}{\texttt{lean/Germain/Naf4073.lean}}
du même dépôt : la formule du feuillet dans tout corps où $2 \neq 0$,
l'identité \eqref{eq:note} dans les entiers, la
conséquence par Nicomaque, le pas $T_{(n+1)^2} - T_{n^2} = n^3 + (n+1)^3$
et les constats sur $189$. L'intégration continue le
compile en traitant tout avertissement comme une erreur et refuse tout axiome
autre que les trois axiomes standard de Lean. La forme démontrée,
$T_n^2 + T_{n+1}^2 = T_{(n+1)^2}$ pour tout $n$, évite la soustraction
tronquée des entiers naturels. Aucune preuve ne dit rien, évidemment, de la
main ni de la date.

\section*{Remerciements}

Je remercie Georg Essl, qui a reconnu dans le manuscrit publié par
Alexanderson l'identité du feuillet de NAF 4073 et me l'a signalé.

\begin{thebibliography}{9}

\bibitem{alexanderson}
G.~L. Alexanderson, «~About the cover: Sophie Germain and a problem in number
theory~», \emph{Bulletin of the American Mathematical Society} 49, n\textsuperscript{o}~2 (avril 2012),
p.~327--331. \href{https://doi.org/10.1090/S0273-0979-2012-01370-8}{doi:10.1090/S0273-0979-2012-01370-8}.

\bibitem{delcentina}
A.~Del~Centina, «~Letters of Sophie Germain preserved in Florence~»,
\emph{Historia Mathematica} 32 (2005), p.~60--75.
\href{https://doi.org/10.1016/j.hm.2003.11.001}{doi:10.1016/j.hm.2003.11.001}.

\bibitem{delcentina-fiocca}
A.~Del~Centina et A.~Fiocca, «~On the correspondence of Sophie Germain~»,
dans M.~T.~Borgato, E.~Neuenschwander et I.~Passeron (dir.),
\emph{Mathematical Correspondences and Critical Editions}, Cham, Springer,
coll. «~Trends in the History of Science~», 2018, p.~147--166.
\href{https://doi.org/10.1007/978-3-319-73577-1_9}{doi:10.1007/978-3-319-73577-1\_9}.

\bibitem{delisle}
L.~Delisle, \emph{Catalogue des manuscrits des fonds Libri et Barrois},
Paris, H.~Champion, 1888.

\bibitem{essl}
G.~Essl, réponse à «~Who wrote this note on triangular numbers in Sophie
Germain's correspondence (BnF NAF 4073)?~», History of Science and
Mathematics, 25 octobre 2026, \url{https://hsm.stackexchange.com/a/19744}.

\bibitem{gallica}
Bibliothèque nationale de France, Nouvelles acquisitions françaises 4073,
\url{https://gallica.bnf.fr/ark:/12148/btv1b10089787b}.

\bibitem{hayes}
B.~Hayes, «~Gauss's day of reckoning~», \emph{American Scientist} 94,
n\textsuperscript{o}~3 (2006), p.~200.
\href{https://doi.org/10.1511/2006.59.200}{doi:10.1511/2006.59.200}.

\bibitem{henry}
C.~Henry, «~Les manuscrits de Sophie Germain et leur récent éditeur.
Documents nouveaux~», \emph{Revue philosophique de la France et de
l'étranger} 8 (1879), p.~619--641.

\bibitem{laubenbacher-pengelley}
R.~Laubenbacher et D.~Pengelley, «~“Voici ce que j'ai trouvé:” Sophie
Germain's grand plan to prove Fermat's Last Theorem~», \emph{Historia
Mathematica} 37 (2010), p.~641--692.
\href{https://doi.org/10.1016/j.hm.2009.12.002}{doi:10.1016/j.hm.2009.12.002}.

\bibitem{oeis}
OEIS Foundation, \emph{The On-Line Encyclopedia of Integer Sequences},
suite A037270, «~$a(n) = n^2(n^2 + 1)/2$~», \url{https://oeis.org/A037270},
consultée le 25 novembre 2026.

\bibitem{omont}
H.~Omont, \emph{Catalogue général des manuscrits français. Nouvelles
acquisitions françaises}, t.~II, \emph{n\textsuperscript{os}~3061--6500}, Paris, E.~Leroux, 1900.

\end{thebibliography}
\end{document}
